\chapter{Especificaciones tecnicas}
\thispagestyle{memoria}\pagestyle{memoria}

\section{Disposiciones generales}
Todos los materiales y equipos deben ser nuevos, de marcas establecidas y
cumplir las reglas del CNE-U y la EM.010. Los valores de catalogo empleados
en este anteproyecto son referenciales; se sustituyen por las placas reales
antes de la construccion.

\section{Conductores y canalizaciones}
Conductores de cobre, aislamiento XLPE o EPR, tension nominal 450/750~V o
superior. Seccion minima 2,5~mm\textsuperscript{2} para fuerza y alumbrado
(Regla 030-002). Las canalizaciones se instalan en tubos PVC SAP o metalicos
segun el metodo de la Tabla 2 (B1 empotrado en pared, D enterrado) y la Regla
070 del CNE-U. Los conductores de proteccion (PE) cumplen la Tabla 16; en
circuitos de motor se adopta igual seccion que la fase.

\section{Cajas, sellos y uniones}
Las cajas deben permitir la instalacion y extraccion de conductores y no
presentar aristas cortantes. En areas clasificadas se aplican sellos y
accesorios compatibles con la zona segun la Seccion 110/120 del CNE-U. Todo
empalme queda dentro de caja accesible.

\section{Tableros y protecciones}
Tableros metalicos con barra de neutro y de tierra separadas (TN-S),
interruptores automaticos de curva C y D segun el cuadro de circuitos, y
proteccion diferencial de 30~mA en circuitos de tomacorrientes y equipos de
area clasificada. El interruptor general de \IPrincipal~y los subalimentadores
se seleccionan de modo que \(I_b \leq I_n \leq I_z\) (Regla 050-104).

\section{Motores y cargas especiales}
Las bombas sumergibles de tanque se alimentan desde el tablero de emergencia
con arranque directo secuencial. Se verifica la coordinacion de arranque con
los subalimentadores (curva D). Los surtidores y sus cabezas electronicas se
alimentan mediante UPS de combustible/control.

\section{Bombas, surtidores y paro de emergencia}
Se incluye circuito de pulsador y paro de emergencia general de playa. El
paro debe desenergizar bombas sumergibles y surtidores en emergencia conforme
a la normativa sectorial de establecimientos de venta al publico de
combustibles y la Regla 120 del CNE-U.

\section{Alumbrado y dispositivos}
Alumbrado LED interior y exterior, incluyendo marquesina, area de despacho y
postes de patio, con circuitos independientes. Alumbrado de emergencia y
senales de evacuacion conforme al articulo 11.1 de la EM.010.

\section{Puesta a tierra, rayo y sobretensiones}
Pozo de puesta a tierra con objetivo de 10~ohm (limite de 25~ohm de la Regla
060-712 del CNE-U), enlace equipotencial de masas, malla de tierra y pararrayo
de 12~m con radio de proteccion de 20~m observado en el plano. Se protege el
tablero general contra sobretensiones transitorias. Los valores quedan sujetos
a resistividad del terreno y medicion con telurimetro.

\section{Grupo electrogeno, ATS y UPS}
Grupo \GrupoModelo~de~\GrupoNom~kVA standby con correccion por altitud,
transferencia automatica de cuatro polos con neutro conmutado, y UPS
senoidales para combustible/control y POS-CCTV-comunicaciones con autonomia
minima de 15 minutos.

\section{Identificacion y documentacion}
Todos los tableros, circuitos, canalizaciones y equipos se identifican con
etiquetas permanentes que coinciden con el cuadro de cargas y los planos. Se
entregara la documentacion de prueba exigida por la EM.010.
